% Part: normal-modal-logic
% Chapter: frame-definability
% Section: definability
%READERNOTE{SOL6-A151: D च्या proof मध्ये V हा model tuple मध्ये आहे; त्याची विशेष निवड करण्याची गरज नाही, असा independence चा अचूक अर्थ घेतला. B आणि 5 च्या proofs मध्ये वापरलेली पद्धत प्रतिपरिवर्तन आहे.}

\documentclass[../../../include/open-logic-section]{subfiles}

\begin{document}

\olfileid{nml}{frd}{def}

\olsection{चौकटींची परिभाष्यता}

\olref[acc]{thm:soundschemas} मधील व्यत्यास प्रतिमानांसाठी
असत्य असले, तरी ``प्रतिमान'' याऐवजी ``चौकट'' घेतल्यास
ते खरे ठरतात: \olref[acc]{thm:soundschemas} मध्ये पाहिलेल्या
गुणधर्मांसाठी, !!a{formula} एखाद्या \emph{चौकटीत} वैध
असेल, तर त्या चौकटीच्या प्राप्यता संबंधाला संबंधित गुणधर्म
\emph{असतोच}. म्हणून संबंधित गुणधर्म असलेल्या चौकटींचे
वर्ग हे !!{formula}s \emph{परिभाषित करतात}.

\begin{defn}
  $\mClass{F}$ हा चौकटींचा वर्ग असेल, तर नेमक्या
  त्यातील चौकटी~$\mModel{F} \in \mClass{F}$
  आणि त्यांच्याच बाबतीत $\mModel{F} \Entails !A$ असेल,
  तर $!A$ हा $\mClass{F}$ वर्ग \emph{परिभाषित करतो},
  असे म्हणतो.
\end{defn}

आता चौकटींसाठी परिभाष्यतेचे पूर्ण निष्कर्ष सिद्ध करू.

\begin{thm}\ollabel{thm:fullCorrespondence}
\olref[acc]{tab:five} च्या उजव्या बाजूचे !!{formula}
चौकट~$\mModel{F}$ मध्ये वैध असेल, तर $\mModel{F}$ ला
डाव्या बाजूचा गुणधर्म असतो.
\end{thm}

\begin{proof}
  \begin{enumerate}
  \item \Ax{D} हे $\mModel{F} = \tuple{W, R}$ मध्ये वैध आहे,
    म्हणजे $\mModel{F} \Entails \Box p \lif \Diamond p$,
    असे समजा. $\mModel{F}$ वर आधारित
    $\mModel{M} = \tuple{W, R, V}$ हे प्रतिमान आणि
    $w \in W$ घ्या. $Rwv$ असलेले एखादे $v$ आहे हे
    दाखवायचे आहे. तसे नाही असे समजा. मग~$p$ सहित
    कोणत्याही~$!A$ साठी $\mSat{M}{\Box !A}[w]$ आणि
    $\mSat/{M}{\Diamond !A}[w]$ ही दोन्ही विधाने खरी
    असतात. पण मग $\mSat/{M}{\Box p \lif
      \Diamond p}[w]$, आणि हे $\mModel{F}
    \Entails \Box p \lif \Diamond p$ या गृहीतकाला
    विरोधी आहे.
  \item \Ax{T} हे $\mModel{F}$ मध्ये वैध आहे,
    म्हणजे $\mModel{F} \Entails \Box
    p \lif p$, असे समजा. कोणतेही जग $w \in W$ घ्या;
    $Rww$ हे दाखवायचे आहे. $u \in V(p)$ तेव्हा आणि
    केवळ तेव्हाच, जेव्हा $Rwu$, असे मूल्यांकन
    ठरवा (आणि $p$ पेक्षा वेगळ्या $q$ साठी $V(q)$
    मनमाना ठेवा, उदाहरणार्थ $V(q) = \emptyset$).
    $\mModel{M} = \tuple{W, R, V}$ घ्या. या रचनेनुसार
    $Rwu$ असलेल्या सर्व $u$ साठी $\mSat{M}{p}[u]$;
    म्हणून $\mSat{M}{\Box p}[w]$. पण गृहीतकावरून
    $\Box p \lif p$ हे~$w$ येथे सत्य आहे; म्हणून
    $\mSat{M}{p}[w]$. $V$ च्या व्याख्येनुसार हे
    $Rww$ असेल तेव्हाच शक्य आहे.
  \item प्रतिपरिवर्तनाने सिद्ध करू: $\mModel{F}$ सममित
    नाही असे समजा. मग \Ax{B}, म्हणजे
    $p \lif \Box\Diamond p$, हे
    $\mModel{F}= \tuple{W, R}$ मध्ये वैध नाही हे
    दाखवू. $\mModel{F}$ सममित नसल्यामुळे
    $u$, $v \in W$ अशी जगं आहेत की $Ruv$ आहे,
    पण $Rvu$ नाही. $w \in V(p)$ तेव्हा आणि केवळ
    तेव्हाच, जेव्हा $Rvw$ नाही, असे $V$ ठरवा (इतरत्र
    $V$ मनमाना आहे). $\mModel{M} =\tuple{W, R,
      V}$ घ्या. $V$ च्या व्याख्येनुसार, $Rvw$
    नसलेल्या सर्व $w$ साठी $\mSat{M}{p}[w]$;
    विशेषतः $Rvu$ नसल्यामुळे $\mSat{M}{p}[u]$.
    तसेच $Rvw$ तेव्हा आणि केवळ तेव्हाच, जेव्हा
    $w \notin V(p)$; म्हणून $Rvw$ आणि
    $\mSat{M}{p}[w]$ या दोन्ही अटी पूर्ण करणारे
    कोणतेही~$w$ नाही. त्यामुळे
    $\mSat/{M}{\Diamond p}[v]$. $Ruv$ असल्याने
    $\mSat/{M}{\Box\Diamond p}[u]$ देखील. म्हणून
    $\mSat/{M}{p \lif
      \Box\Diamond p}[u]$, आणि त्यामुळे
    $\Ax{B}$ हे~$\mModel{F}$ मध्ये वैध नाही.
  \item \Ax{4} हे $\mModel{F} = \tuple{W,R}$ मध्ये
    वैध आहे, म्हणजे $\mModel{F} \Entails \Box p
    \lif \Box\Box p$, असे समजा. $u$, $v$, $w \in W$
    ही कोणतीही जगं $Ruv$ आणि $Rvw$ अशा अटी पूर्ण
    करत असू द्या; $Ruw$ सिद्ध करायचे आहे.
    $z \in V(p)$ तेव्हा आणि केवळ तेव्हाच, जेव्हा
    $Ruz$, असे $V$ ठरवा (इतरत्र $V$ मनमाना आहे).
    $\mModel{M} =\tuple{W, R, V}$ घ्या. $V$ च्या
    व्याख्येनुसार $Ruz$ असलेल्या सर्व $z$ साठी
    $\mSat{M}{p}[z]$; म्हणून $\mSat{M}{\Box p}[u]$.
    गृहीतकावरून \Ax{4}, म्हणजे $\Box p
    \lif \Box \Box p$, हे $u$ येथे सत्य आहे;
    म्हणून $\mSat{M}{\Box \Box
      p}[u]$. $Ruv$ आणि $Rvw$ असल्यामुळे
    $\mSat{M}{p}[w]$; पण $V$ च्या व्याख्येनुसार
    हे $Ruw$ असेल तेव्हाच शक्य आहे.
  \item पुन्हा प्रतिपरिवर्तन वापरू. चौकट
    $\mModel{F} = \tuple{W, R}$ युक्लिडी नाही असे
    मानून ती \Ax{5} असत्य ठरवते, म्हणजे
    $\mModel{F} \Entails/ \Diamond p \lif
      \Box\Diamond p$, हे दाखवू. $u$, $v$, $w \in W$
    अशी जगं असू द्या की $Rwu$ आणि $Rwv$ आहेत,
    पण $Ruv$ नाही. सर्व जगं $z$ साठी $z \in V(p)$
    तेव्हा आणि केवळ तेव्हाच, जेव्हा $Ruz$ असे
    \emph{नाही}, असे $V$ ठरवा. $\mModel{M}
    =\tuple{W, R, V}$ घ्या. मग गृहीतकावरून
    $\mSat{M}{p}[v]$, आणि $Rwv$ असल्यामुळे
    $\mSat{M}{\Diamond p}[w]$ देखील. परंतु
    $Ruy$ आणि $\mSat{M}{p}[y]$ असलेले कोणतेही
    जग~$y$ नाही; म्हणून $\mSat/{M}{\Diamond
      p}[u]$. $Rwu$ असल्यामुळे
    $\mSat/{M}{\Box\Diamond
      p}[w]$; म्हणून \Ax{5}, म्हणजे
    $\Diamond p \lif \Box\Diamond p$, हे~$w$
    येथे असत्य आहे.
  \end{enumerate}
\end{proof}

\Ax{D} ची सिद्धता आणि इतर प्रकरणे यांत एक फरक दिसेल:
त्यात मूल्यांकन~$V$ ची कोणतीही विशेष निवड करावी लागली नाही. प्रत्यक्षात,
$\mSat{M}{\Ax{D}}$ असल्यास $\mModel{M}$ सर्वत्र
उत्तरवर्ती आहे, हे आपण सिद्ध केले. म्हणून \Ax{D}
हा रूपबंध केवळ चौकटींचा नव्हे, तर सर्वत्र उत्तरवर्ती
\emph{प्रतिमानांचा} वर्ग परिभाषित करतो.

\begin{cor}\ollabel{prop:D-serial}
  \Ax{D} सत्य असलेले कोणतेही प्रतिमान सर्वत्र उत्तरवर्ती असते.
\end{cor}

\begin{cor}
\olref[acc]{tab:five} च्या उजव्या बाजूचे प्रत्येक
!!{formula} डावीकडील गुणधर्म असलेल्या चौकटींचा
वर्ग परिभाषित करते.
\end{cor}

\begin{proof}
  \olref[acc]{thm:soundschemas} मध्ये आपण सिद्ध केले की,
  प्रतिमानाला डावीकडील गुणधर्म असेल, तर उजवीकडचे
  !!{formula} त्यात सत्य असते. म्हणून चौकट~$\mModel{F}$
  ला तो गुणधर्म असेल, तर उजवीकडचे !!{formula}
  $\mModel{F}$ मध्ये वैध असते.
  \olref{thm:fullCorrespondence} मध्ये आपण व्यत्यास
  सिद्ध केले: उजवीकडचे !!a{formula}~$\mModel{F}$
  मध्ये वैध असेल, तर $\mModel{F}$ ला डावीकडचा
  गुणधर्म असतो.
\end{proof}

\begin{prob}
\olref[nml][frd][acc]{tab:anotherfive} च्या उजव्या
बाजूचे !!{formula} चौकट~$\mModel{F}$ मध्ये वैध असेल,
तर $\mModel{F}$ ला डावीकडील गुणधर्म असतो, हे दाखवा.
यासाठी डावीकडील गुणधर्म \emph{नसलेली} चौकट घ्या,
आणि उजवीकडचे !!{formula} एखाद्या जगात असत्य ठरेल
असे योग्य~$V$ ठरवा.
\end{prob}

\olref{thm:fullCorrespondence} वरून गुणधर्म एकत्रही
करता येतात हे दिसते: उदाहरणार्थ \Ax{B} आणि
\Ax{4} हे दोन्ही~$\mModel{F}$ मध्ये वैध असतील,
तर चौकट सममित आणि संक्रमक दोन्ही असते.
अनेक महत्त्वाची मोडल तर्कशास्त्रे काही चौकट-गुणधर्म
एकत्र असलेल्या सर्व चौकटींमध्ये वैध असलेल्या
!!{formula}s चे संच म्हणून निश्चित होतात. म्हणून
त्यांना संबंधित परिभाषक !!{formula}s वैध असलेल्या
सर्व चौकटींमध्ये वैध असलेल्या !!{formula}s चे संच
म्हणूनही सांगता येते. उदाहरणार्थ, अभिजात व्यवस्था
\Log{S4} म्हणजे सर्व परावर्ती आणि संक्रमक चौकटींमध्ये
वैध असलेल्या !!{formula}s चा संच; म्हणजे \Ax{T}
आणि~\Ax{4} दोन्ही वैध असलेल्या सर्व चौकटींमध्ये
वैध असलेल्या सूत्रांचा संच.
\Log{S5} म्हणजे सर्व परावर्ती, सममित आणि युक्लिडी
चौकटींमध्ये वैध असलेल्या सूत्रांचा संच; म्हणजे
\Ax{T}, \Ax{B}, आणि \Ax{5} ही तिन्ही वैध
असलेल्या सर्व चौकटींमध्ये वैध असलेल्या सूत्रांचा संच.

$R$ च्या गुणधर्मांमधील तार्किक संबंध सामान्यतः
संबंधित परिभाषक !!{formula}s मधील संबंधांशी जुळतात.
उदाहरणार्थ, प्रत्येक परावर्ती संबंध सर्वत्र उत्तरवर्ती
असतो; म्हणून चौकटीत \Ax{T} वैध असेल, तर
\Ax{D} देखील वैध असतो. (हा संबंध \emph{तार्किक
निष्पन्नतेचा नाही}, हे लक्षात घ्या. $\mSat{M}{\Ax{T}}[w]$
असल्यावर नेहमी $\mSat{M}{\Ax{D}}[w]$ असते, असे
नाही.) असे काही संबंध पुढे नोंदवतो.

\begin{prop}\ollabel{prop:relation-facts}
  $R$ हा $W$ या संचावरील द्विस्थानी संबंध असू द्या. मग:
  \begin{enumerate}
  \item $R$ परावर्ती असेल, तर तो सर्वत्र उत्तरवर्ती असतो.
  \item $R$ सममित असेल, तर तो संक्रमक असतो तेव्हा आणि
    केवळ तेव्हाच, जेव्हा तो युक्लिडी असतो.
  \item $R$ सममित किंवा युक्लिडी असेल, तर तो दुर्बल
    दिशित असतो (त्याला ``हिऱ्याचा गुणधर्म'' असतो).
  \item $R$ युक्लिडी असेल, तर तो दुर्बल संयोजित असतो.
  \item $R$ फलनीय असेल, तर तो सर्वत्र उत्तरवर्ती असतो.
  \end{enumerate}
\end{prop}

\begin{prob}
  \olref[nml][frd][def]{prop:relation-facts} सिद्ध करा.
\end{prob}

\end{document}
